% STATUS: rewritten 2026-11-05; every claim is cited.
\section{Near-rings and prime near-rings with derivation}\label{sec:other-systems:near-rings}

The previous section listed the systems of this chapter, and the first of
them is the near-ring \cite[p.~31]{BellMason1987}. A ring has two
distributive laws; a near-ring keeps only one
\cite[p.~31]{BellMason1987}. The product rule of a derivation still makes
sense, but the proofs from prime rings that use the missing law break
\cite[p.~31]{BellMason1987}. Bell and Mason asked in 1987
what survives \cite[p.~31]{BellMason1987}. Their setting follows.

\begin{definition}[{\cite[p.~31]{BellMason1987}}]\label{def:other:near-ring-derivation}
Let $N$ be a zero-symmetric left near-ring with multiplicative center $Z$.
A \emph{derivation} on $N$ is an additive endomorphism $D$ of $(N, +)$
satisfying the product rule $D(xy) = x\,D(y) + D(x)\,y$ for all $x, y \in N$.
An element $x \in N$ with $D(x) = 0$ is called a \emph{constant}. The
derivation $D$ is \emph{commuting} if $[x, D(x)] = 0$ for all $x \in N$.
\end{definition}

The commutativity theorems below need the near-ring analogue of a prime
ring \cite[p.~31]{BellMason1987}.

\begin{definition}[{\cite[p.~31]{BellMason1987}}]\label{def:other:prime-near-ring}
The near-ring $N$ is \emph{prime} if $a\,N\,b = \{0\}$ implies $a = 0$ or
$b = 0$ for all $a, b \in N$.
\end{definition}

Bell and Mason note that this condition implies the usual definition of a
prime near-ring \cite[p.~31]{BellMason1987}.

The first surprise is that the product rule alone gives a weak form of the
missing distributive law.

\begin{lemma}[{\cite[Lemma~1, p.~32]{BellMason1987}}]\label{lem:other:partial-distributive}
Let $D$ be a derivation on a near-ring $N$. Then $N$ satisfies the partial
distributive law
\[
\bigl(a\,D(b) + D(a)\,b\bigr)\,c = a\,D(b)\,c + D(a)\,b\,c
\qquad (a, b, c \in N).
\]
\end{lemma}

Bell and Mason proved the first commutativity theorem for a near-ring with
a derivation \cite[Theorem~1, p.~32]{BellMason1987}.

\begin{theorem}[{\cite[Theorem~1, p.~32]{BellMason1987}}]\label{thm:other:bell-mason-1}
Let $N$ have no nonzero divisors of zero. If $N$ admits a nontrivial
commuting derivation $D$, then $(N, +)$ is abelian.
\end{theorem}

They then moved to prime near-rings, where the analogue of the first
Posner theorem on rings holds \cite[Theorem~2, p.~33]{BellMason1987}.

\begin{theorem}[{\cite[Theorem~2, p.~33]{BellMason1987}}]\label{thm:other:bell-mason-2}
If a prime near-ring $N$ admits a nontrivial derivation $D$ for which
$D(N) \subseteq Z$, then $(N, +)$ is abelian. Moreover, if $N$ is
$2$-torsion-free, then $N$ is a commutative ring.
\end{theorem}

The same conclusion follows when the derivations commute with each other
\cite[Theorem~3, p.~33]{BellMason1987}.

\begin{theorem}[{\cite[Theorem~3, p.~33]{BellMason1987}}]\label{thm:other:bell-mason-3}
Let $N$ be a prime near-ring admitting a nontrivial derivation $D$ such
that $[D(x), D(y)] = 0$ for all $x, y \in N$. Then $(N, +)$ is abelian.
If $N$ is $2$-torsion-free as well, then $N$ is a commutative ring.
\end{theorem}

Bell continued alone in 1997 \cite[p.~191]{Bell1997II}. His condition now
holds only on a subset \cite[p.~191]{Bell1997II}. He called a nonempty
subset $U$ of $N$ a semigroup
right ideal if $U N \subseteq U$, a semigroup left ideal if $N U \subseteq U$,
and a semigroup ideal if both hold \cite[p.~191]{Bell1997II}.

\begin{proposition}[{\cite[Abstract, p.~191]{Bell1997II}}]\label{prop:other:bell-ii}
Let $N$ be a $3$-prime zero-symmetric left near-ring, and let $d$ be a
nonzero derivation on $N$. Let $U$ be a nonzero subset of $N$. Then $N$ is
a commutative ring if one of the following holds: \textup{(a)} $U$ is a
semigroup right or left ideal of $N$ and $d(U)$ is multiplicatively
central; \textup{(b)} $U$ is a semigroup ideal of $N$, $d^{2} \neq 0$ and
$[d(U), d(U)] = \{0\}$.
\end{proposition}

Arga\c{c} asked in 1997 when a derivation on a near-ring is a homomorphism
on a subset \cite[p.~737]{Argac1997}. For rings, Bell and Kappe had shown
that such a derivation
vanishes \cite[p.~737]{Argac1997}. The same holds for near-rings
\cite[p.~737]{Argac1997}.

\begin{theorem}[{\cite[Theorem, p.~737]{Argac1997}}]\label{thm:other:argac}
Let $N$ be a semiprime right near-ring, and let $d$ be a derivation on $N$.
Let $A$ be a subset of $N$ such that $0 \in A$ and $A N \subseteq A$. If $d$
acts as a homomorphism on $A$ or as an anti-homomorphism on $A$, then
$d(A) = \{0\}$.
\end{theorem}

In a near-ring the prime ideals behind these theorems come in four
strengths, listed by Groenewald in 1991
\cite[Definition~1, p.~2669]{Groenewald1991}.

\begin{definition}[{\cite[Definition~1, p.~2669]{Groenewald1991}}]\label{def:other:prime-ideals}
Let $N$ be a near-ring and $P$ an ideal of $N$. Then $P$ is
\emph{$0$-prime} if $A B \subseteq P$ implies $A \subseteq P$ or
$B \subseteq P$ for all ideals $A, B$ of $N$; \emph{$1$-prime} if the same
holds for all left ideals $A, B$; \emph{$2$-prime} if the same holds for
all $N$-subgroups $A, B$; and \emph{$3$-prime} if $a N b \subseteq P$
implies $a \in P$ or $b \in P$ for all $a, b \in N$. The near-ring $N$ is
\emph{$\nu$-prime} ($\nu = 0, 1, 2, 3$) if the zero ideal is $\nu$-prime.
\end{definition}

The $3$-prime near-rings are the ones named 3-prime in the 1997 paper of
Bell above \cite[p.~191]{Bell1997II}. The next section moves to
$\Gamma$-rings, where the product itself becomes a family
\cite[p.~2]{DeyPaulRakhimov2012}.

A multiplicative semiderivation is introduced on a semigroup ideal of a 3-prime near-ring, and the commutativity conditions it forces there are proved~\cite{WlSemigroupIdealsWith2022}.

A 3-prime near-ring is commutative when a multiplicative generalized derivation of it satisfies any of five stated conditions on images and commutators~\cite{WlOnCommutativityOf2019}.

The commutativity results for Jordan ideals of a 3-prime near-ring are extended from derivations to left derivations, which the author shows is enough~\cite{WlOnJordanIdeals2023}.

A 3-prime near-ring is commutative when a multiplicative generalized derivation of it satisfies any of five stated conditions on images and commutators~\cite{WlOnCommutativityOf2019}.

The commutativity results for Jordan ideals of a 3-prime near-ring are extended from derivations to left derivations, which the author shows is enough~\cite{WlOnJordanIdeals2023}.
